#HTML Learning notes

## Normal elements vs. Reemplace elements

- **Normal elements:** They have an opening tag and a closing tag. For example, `<p>...</p>`, `<div>...</div>`, etc.
- **Reemplace elements:** They do not have a closing tag and are **self-closing**. For example, `<img />`, `<br />`, etc.


## Meta tags
- **Meta tags** are used to provide metadata about the HTML document. They are placed inside the `<head>` section of the HTML document and are not displayed on the page but can be used by browsers, search engines, and other web services.

### SEO (Search Engine Optimization) meta tags
**SEO meta tags** are used to improve the visibility of a webpage in search engine results. They provide information about the content of the page, which can help search engines understand and rank the page appropriately.
- Examples of SEO meta tags include:
    - `<meta name="description" content="A brief description of the page's content.">`
    - `<meta name="keywords" content="HTML, CSS, JavaScript, web development">`
    - `<meta name="author" content="Bl4ackM0oN">`
    - `<meta name="viewport" content="width=device-width, initial-scale=1.0">`
    - `<meta name="robots" content="index, follow">`

### Open Graph meta tags
**Open Graph meta tags** are used to control how a webpage is displayed when shared on social media platforms. They provide information about the title, description, image, and other attributes of the page that can be used by social media platforms to create rich previews.
- Examples of Open Graph meta tags include:
    - `<meta property="og:title" content="Bl4ackM0oN Portfolio">`
    - `<meta property="og:description" content="Bl4ackM0oN's portfolio webpage: showcasing experience and passion for learning, creating, and innovating.">`
    - `<meta property="og:image" content="./images/Myself.png">`
    - `<meta property="og:image-alt" content="Bl4ackM0oN's profile picture">`

### Link elements
**Link elements** are used to link external resources to the HTML document. They can be used for various purposes, such as linking stylesheets, icons, alternate versions of the page, and canonical versions of the page.
- Examples of link elements include:
    - `<link rel="icon" href="./images/Myself.png" type="image/png">`
    - `<link rel="alternate" href="" hreflang="es">`
    - `<link rel="canonical" href="">`
    - `<link rel="stylesheet" href="styles.css">` -> This is used to link an external CSS stylesheet to the HTML document.


## Semantic HTML
**Semantic HTML** refers to the use of HTML elements that convey meaning about the content they contain. These elements help both developers and browsers understand the structure and purpose of the content, improving accessibility and SEO.

### '<span>' vs. '<div>' vs. '<section>'
- `<span>` is an inline element used to group and style a **portion of text** or other inline elements without affecting the layout.
- `<div>` is a block-level element used to group and style **larger sections of content**, such as paragraphs, images, or other block-level elements.
- `<section>` is a semantic block-level element used to group **related content together**, often with a heading, to indicate that the content belongs to a specific section of the page.

### '<strong>' vs. '<b>'
- `<strong>` is used to indicate that the text is of strong **importance**.
- `<b>` is used to simply make the text bold without conveying any additional meaning or importance.

### '<em>' vs. '<i>'
- `<em>` is used to indicate that the text has **emphasis**.
- `<i>` is used to simply make the text italic without conveying any additional meaning or importance.

### role attribute
The `role` attribute is used to define the **purpose** of an element in terms of accessibility. It helps assistive technologies understand the function of an element, which can improve the user experience for people with disabilities. For example, `role="navigation"` indicates that the element is a navigation landmark, while `role="button"` indicates that the element functions as a button.

## Interesting elements
### details and summary
The `<details>` element is used to create a disclosure widget that can be toggled open or closed by the user. 
The `<summary>` element is used to provide a summary or caption for the content inside the `<details>` element. When the user clicks on the summary, the details content will be revealed or hidden.

### video
The `<video>` element is used to embed video content in an HTML document. 
It provides a standard way to include video files and offers various attributes for controlling playback, such as `controls`, `autoplay`, `loop`, and `muted`. 

### audio
The `<audio>` element is used to embed audio content in an HTML document.
It provides a standard way to include audio files and offers various attributes for controlling playback.

##iframe
The `<iframe>` element is used to embed another HTML document within the current document. 
It allows you to display content from another source, such as a different webpage, a video, or a map, without leaving the current page. The `src` attribute is used to specify the URL of the content to be embedded.

## dialog
The `<dialog>` element is used to create a dialog box or modal window in an HTML document.
It provides a standard way to create interactive dialogs that can be opened and closed by the user.

## Interesting attributes
### loading='lazy'
The `loading="lazy"` attribute is used on `<img>` and `<iframe>` elements to enable lazy loading. 
This means that the browser will defer loading the image or iframe until it is about to enter the viewport, which can improve page load times and reduce bandwidth usage for users who do not scroll all the way down the page.